\documentclass[11pt,a4paper]{article}
\usepackage[T1]{fontenc}
\usepackage[utf8]{inputenc}
\usepackage[main=italian,provide=*]{babel}
\usepackage{amsmath,amssymb}
\usepackage{geometry}
\geometry{margin=2.3cm,top=2.5cm,bottom=2.7cm}
\usepackage{booktabs}
\usepackage[table,dvipsnames]{xcolor}
\usepackage{tcolorbox}
\tcbuselibrary{skins,breakable}
\usepackage{titlesec}
\usepackage{enumitem}
\usepackage{needspace}
\usepackage{fancyhdr}
\usepackage{graphicx}
\usepackage{hyperref}
\hypersetup{colorlinks=true,linkcolor=Sepia,citecolor=Sepia,urlcolor=Sepia}
\newcommand{\Tr}{\operatorname{Tr}}
\newcommand{\ket}[1]{\lvert #1 \rangle}
\newcommand{\bra}[1]{\langle #1 \rvert}

\definecolor{inkblue}{HTML}{1B3A5C}
\definecolor{lightblue}{HTML}{EAF1F8}
\definecolor{accent}{HTML}{B0562A}
\definecolor{softgray}{HTML}{6B6B6B}
\definecolor{okgreen}{HTML}{2E6B3E}
\definecolor{okgreenbg}{HTML}{EAF5EC}

\titleformat{\section}
  {\normalfont\Large\bfseries\color{inkblue}}
  {\thesection}{0.6em}{}[{\color{inkblue!40}\titlerule}]
\titlespacing{\section}{0pt}{0.9em}{0.4em}
\titleformat{\subsection}
  {\normalfont\large\bfseries\color{accent}}
  {}{0em}{}
\titlespacing{\subsection}{0pt}{0.6em}{0.2em}

\pagestyle{fancy}
\fancyhf{}
\renewcommand{\headrulewidth}{0.4pt}
\fancyhead[L]{\small\color{softgray}VQE con termine DM sotto rumore (anello)}
\fancyhead[R]{\small\color{softgray}\thepage}
\fancyfoot[C]{\small\color{softgray}Tesi triennale, Samuele --- Relatore: Prof.\ Alessandro Chiesa}

\newtcolorbox{keyeq}[1][]{
  colback=lightblue, colframe=inkblue!70, boxrule=0.6pt,
  arc=2pt, left=8pt, right=8pt, top=3pt, bottom=3pt, #1
}
\newtcolorbox{okbox}[1]{
  colback=okgreenbg, colframe=okgreen!60, boxrule=0.6pt,
  arc=2pt, left=7pt, right=7pt, top=2pt, bottom=2pt,
  fonttitle=\bfseries\color{okgreen}, title={#1}
}

\title{\vspace{-1.5em}\color{inkblue}VQE con termine DM sotto rumore --- trimero ad anello\\[0.2em]
\Large\normalfont\color{softgray}Quanto degrada un risultato ideale su un simulatore rumoroso, a 3 qubit}
\author{Note di lavoro --- Tesi triennale, Samuele\\ Relatore: Prof.\ Alessandro Chiesa}
\date{}

\begin{document}
\sloppy
\maketitle
\thispagestyle{fancy}
\vspace{-2.0em}

File di codice: \texttt{vqe\_dm\_\allowbreak rumoroso\_trimero\_anello.py}
(mirror di \texttt{vqe\_dm\_rumoroso\_dimero.py}). Punto di lavoro:
$J=1,\,J'=0.4,\,b=b_c=2.4,\,D=0.15$ (mode ``B'', confermato in Parte 1).
Modello di rumore: \texttt{noise\_model\_trimero\_anello.py} (wrapper su
\texttt{noise\_model\_dimero.py} --- Stadio 1), ai valori di riferimento
\texttt{ibm\_torino} salvo dove diversamente indicato.

\tableofcontents

\section{Obiettivo}

Stessa domanda dello Stadio 2 sul dimero, qui sul sistema a 3 spin: il
VQE ideale (Parte 1, nessun rumore) trova parametri
$\theta^*_\text{ideale}$ che preparano lo stato fondamentale
dell'Hamiltoniana dell'anello con DM entro la precisione numerica ---
fedeltà $0.9999999999996$, registrata in \texttt{w2q6\_params\_
optimal.npz}. Se lo stesso circuito, con la stessa ricetta di parametri,
viene eseguito su un simulatore \textbf{rumoroso}, quanto degradano
energia e fedeltà?

\begin{keyeq}[title={Scelta di modellazione dichiarata, identica al dimero}]
\centering
Si riusano i parametri VQE già ottimizzati nel caso ideale, \textbf{senza}
rioptimizzare sotto rumore: questo documento misura la degradazione di un
risultato ideale, non la convergenza di un VQE noise-aware (estensione
facoltativa, non nello scope attuale per il trimero --- vedi
\texttt{domande\_relatore.md}).
\end{keyeq}

\section{Osservabili: energia e fedeltà rumorose}

Stessa formalizzazione del dimero, qui su un registro a 3 qubit:
$E_\text{rumoroso}=\Tr[\rho H]$,
\begin{keyeq}
\centering
$F(\rho,\ket{\psi_\text{ideale}}) = \bra{\psi_\text{ideale}}\rho\ket{\psi_\text{ideale}}$
\end{keyeq}
La derivazione della semplificazione (fedeltà di Uhlmann per un
riferimento puro) non è ripetuta qui: è già in
\texttt{risultati\_vqe\_dm\_\allowbreak rumoroso\_dimero.tex}. Qui si
verifica solo che valga anche a dimensione $d=8$ (3 qubit), non solo alla
$d=4$ del dimero (Sez.~\ref{sec:verifica}).

\section{Risultati}

\begin{center}
\renewcommand{\arraystretch}{1.3}
\begin{tabular}{@{}lccc@{}}
\toprule
\textbf{Condizione} & $\boldsymbol{E}$ & $\boldsymbol{F}$ & \textbf{CNOT}\\
\midrule
Rumore nullo (verifica di coerenza) & $-5.53437002$ & $1.0000000000$ & $6$\\
Rumore di riferimento (\texttt{ibm\_torino}) & $-5.37687489$ & $0.9702147050$ & $6$\\
\bottomrule
\end{tabular}
\end{center}

\begin{okbox}{Verificato}
Il caso a rumore nullo riproduce $E_0=-5.53437002$ (l'energia esatta,
Parte 1) a precisione numerica --- coerenza interna della pipeline. Al
rumore di riferimento: scostamento di energia $+0.157495$, perdita di
fedeltà $1-F=0.029785$ ($\sim3.0\%$) --- circa il \textbf{doppio} della
perdita osservata sul dimero ($1-F=0.014994$, $\sim1.5\%$), coerente con
un ansatz a 6 CNOT invece di 2: più gate a 2 qubit, più occasioni per il
canale depolarizzante di agire.
\end{okbox}

\section{Dipendenza da \texorpdfstring{$\varepsilon_{2q}$}{eps\_2q}}
\label{sec:verifica}

\begin{figure}[h]
\centering
\includegraphics[width=0.92\linewidth]{figure/fig_scan_eps2q_trimero_anello.pdf}
\caption{\textit{Energia e fedeltà rumorose in funzione di
$\varepsilon_{2q}$ (a $\varepsilon_{1q}=2.9\times10^{-4}$,
$p_\text{readout}=2.3\times10^{-2}$ fissati). La linea tratteggiata
verticale segna il valore di riferimento
$\varepsilon_{2q}=3.8\times10^{-3}$. Degradazione monotona e liscia,
nessuna soglia o comportamento anomalo nell'intervallo esplorato ---
stesso andamento qualitativo del dimero, pendenza diversa (piu' ripida,
coerente col maggior numero di CNOT).}}
\end{figure}

\begin{okbox}{Verificato (formula della fedeltà a $d=8$)}
$F=\bra{\psi}\rho\ket{\psi}$ confrontata contro
\texttt{qiskit.quantum\_info.state\_fidelity} su uno stato misto casuale
a $d=8$ (non il caso specifico del VQE, per evitare una verifica
circolare, e non la stessa dimensione $d=4$ già verificata sul dimero):
differenza $0$ a precisione di macchina.
\end{okbox}

\section{Cosa non è ancora stato controllato qui}

Identico al dimero: nessuna dinamica qui (solo lo stato fondamentale
statico), nessuna rioptimizzazione noise-aware. Serve solo come
riferimento di degradazione per uno stato preparato, prima di introdurre
l'evoluzione temporale (Stadio 3).

\section{Considerazioni e cosa serve al passo successivo}

Due fatti usati esplicitamente allo Stadio 3 (quantum simulation Trotter
sotto rumore): (1) i parametri $\theta^*_\text{ideale}$ e il loro
CNOT-count minimo (6, per la sola preparazione) sono il punto di partenza
--- lo stadio successivo aggiunge $N$ passi di Trotter sopra questa
stessa preparazione; (2) è naturale chiedersi se la degradazione qui
osservata ($1-F\sim3.0\%$ al rumore di riferimento) entri nella fedeltà
totale di Trotter come contributo \emph{additivo e costante} in $N$ ---
ma questa è un'\textbf{ipotesi da verificare}, non un fatto già stabilito:
sul dimero la stessa decomposizione additiva era stata verificata per un
\emph{correlatore} (Stadio 4), poi estesa per analogia alla fedeltà di
Trotter (Stadio 3) senza una derivazione indipendente per quella metrica
--- un errore di ragionamento poi corretto esplicitamente
(\texttt{log\_decisioni.md}, sessione 5 settembre). Lo Stadio 3 di questo
documento non deve ripetere lo stesso errore.

\begin{okbox}{Sintesi in una riga}
Riusare parametri VQE ideali sotto rumore di riferimento costa
$\sim3.0\%$ di fedeltà con un ansatz a 6 CNOT --- circa il doppio del
dimero (2 CNOT, $\sim1.5\%$), coerente col maggior conteggio di gate a 2
qubit.
\end{okbox}

\needspace{4\baselineskip}
\section*{File prodotti}
\addcontentsline{toc}{section}{File prodotti}

\texttt{vqe\_dm\_\allowbreak rumoroso\_trimero\_anello.py},
\texttt{validate\_vqe\_dm\_\allowbreak rumoroso\_trimero\_anello.py},
\texttt{generate\_figure\_\allowbreak scan\_eps2q\_trimero\_anello.py},
questo documento.

\end{document}
